\documentclass[11pt]{article}
\usepackage[T1]{fontenc}
\usepackage[utf8]{inputenc}
\usepackage[margin=1in]{geometry}
\usepackage{hyperref}
\title{Research report: Does grounding answers in peer-reviewed literature improve accuracy?}
\author{Autonomous Multimodal Scientific Discovery Agent}
\date{2026-11-04}
\begin{document}
\maketitle
\begin{abstract}
Does grounding answers in peer-reviewed literature improve accuracy? Statistically significant results were obtained for: Accuracy: grounded vs ungrounded. No statistically significant result was found for: Latency: grounded vs ungrounded (seconds). This does not show that there is no effect. Significance alone does not establish that the hypothesis is correct; the direction and size of each effect should be checked against it.
\end{abstract}
\section{Objective}
Research question: Does grounding answers in peer-reviewed literature improve accuracy?

Hypothesis: Grounding answers in peer-reviewed sources raises answer accuracy \cite{src1}\cite{src2}.

Evidence base: 2 cited source(s).

\section{Methods}
\begin{itemize}
\item Experiment: Eight runs per condition on the same question set.
\item Execution: completed successfully after 2 attempt(s).
\item Statistical analysis: Welch two-sample t-test (two-sided, alpha = 0.05).
\item Holm-Bonferroni correction was applied across 2 analyses.
\item Sample sizes: Accuracy: grounded vs ungrounded: grounded = 8, ungrounded = 8; Latency: grounded vs ungrounded (seconds): grounded = 6, ungrounded = 6.
\end{itemize}
\section{Findings}
\begin{itemize}
\item Accuracy: grounded vs ungrounded: Welch two-sample t-test, t = 7.64, p < 0.0001, Holm-adjusted p < 0.0001, Cohen's d = 3.8188 (large), 95\% CI [0.036, 0.064] for mean(grounded) - mean(ungrounded). grounded has the higher mean (difference 0.05). Statistically significant at alpha = 0.05.
\item Latency: grounded vs ungrounded (seconds): Welch two-sample t-test, t = 0.784, p = 0.4510, Holm-adjusted p = 0.4510, Cohen's d = 0.4529 (small), 95\% CI [-0.1227, 0.256] for mean(grounded) - mean(ungrounded). grounded has the higher mean (difference 0.06667). Not statistically significant at alpha = 0.05; this does not show that there is no effect.
\item Automated figure analysis, Figure 3 of \cite{src2}: hallucination rate falls as retrieval depth increases
\end{itemize}
\section{Conclusion}
Statistically significant results were obtained for: Accuracy: grounded vs ungrounded. No statistically significant result was found for: Latency: grounded vs ungrounded (seconds). This does not show that there is no effect. Significance alone does not establish that the hypothesis is correct; the direction and size of each effect should be checked against it.

\section{Limitations}
\begin{itemize}
\item A p-value is not the probability that the hypothesis is true, and p below alpha does not guarantee the effect is real or important.
\item Small sample (smallest group has 6 values): low power and unstable estimates. Treat this result as preliminary.
\item This report was generated automatically from the experiment outputs and retrieved literature and should be reviewed by a person before it is relied on.
\end{itemize}
\begin{thebibliography}{99}
\bibitem{src1} Lewis, Perez, Piktus et al. (2020). Retrieval-augmented generation for knowledge-intensive NLP tasks. NeurIPS. \url{https://doi.org/10.5555/3495724.3496517}
\bibitem{src2} Ji, Lee (2023). Survey of hallucination in natural language generation. ACM Computing Surveys. \url{https://doi.org/10.1145/3571730}
\end{thebibliography}
\end{document}
